\section{Admissible translates for \texorpdfstring{$B\oplus B+t$}{B⊕B+t} in abelian groups}
\label{G:sec:translates}

In this section we answer Question~5.16 of~\cite{kra2023problems} and the related
Question~\S5.3(a). For a countably infinite abelian group $G$ in which $2G=\{g+g:g\in G\}$
has finite index (by work of Charamaras--Mountakis and Ackelsberg these are exactly the
groups in which Conjecture~5.15 of~\cite{kra2023problems} holds), we show that a set
$T\subset G$ of translates suffices if and only if $T$ meets every coset of $2G$. In
particular, the transversals of $2G$ are exactly the minimal sets of translates that
suffice. We also determine exactly which sets arise as the set of admissible translates
of some $A$ of positive upper Banach density, and we give two-sided bounds for the set of
admissible translates of each individual $A$.
The new input is small: the existence of some translate is a theorem of
Charamaras--Mountakis~\cite{charamaras2024finding}, and the two elementary observations
we use have already appeared in the literature in closely related settings; see
\cref{G:rem:credit}.

\subsection{Setting and the survey's questions}
Let $(G,+)$ be a countable abelian group. A \emph{\Fol{} sequence} in $G$ is a sequence
$\Phi=(\Phi_N)_{N\in\N}$ of finite nonempty subsets of $G$ such that
$|\Phi_N\cap(g+\Phi_N)|/|\Phi_N|\to1$ for every $g\in G$. (Since $G$ is abelian, left,
right and two-sided \Fol{} sequences coincide.) Every countable abelian group is amenable,
so \Fol{} sequences exist. For $A\subset G$ we write
$\udens_\Phi(A)=\limsup_{N\to\infty}|A\cap\Phi_N|/|\Phi_N|$, and we define the
\emph{upper Banach density} as
\[
\ubd(A)=\sup_\Phi\udens_\Phi(A),
\]
the supremum being taken over all \Fol{} sequences in $G$. Thus $\ubd(A)>0$ if and only if
$\udens_\Phi(A)>0$ for some \Fol{} sequence $\Phi$; this is the notion of positive upper
Banach density used in~\cite{charamaras2024finding}. For $B\subset G$ let
$B\oplus B=\{b+b':b,b'\in B,\ b\neq b'\}$.

Conjecture~5.15 of~\cite{kra2023problems} asserts that in every countable abelian group
$G$, every set $A\subset G$ with $\ubd(A)>0$ contains $B\oplus B+t$ for some infinite
$B\subset G$ and some $t\in G$. The survey then asks:

\begin{question}{5.16}\label{G:q:516}
Assuming that Conjecture~5.15 holds, what values of the shift $t$ suffice?
\end{question}

\begin{question}{\S5.3(a)}\label{G:q:53a}
The authors write: ``We suspect that there is a direct relation between coset
representatives of the subgroup $\{g+g:g\in G\}$ and the restrictions that can be placed
on the shift $t$.'' Make this relation precise.
\end{question}

The sentence ``We would like to know if similar restrictions can be placed in other
groups'', which precedes Question~5.16 in~\cite{kra2023problems}, is the motivation for
Question~5.16 itself. By Question~\S5.3(a) we mean the suspicion quoted above, which the
survey states right after Question~5.16.

As recorded in~\cite{kra2023problems}, Conjecture~5.15 is known to hold precisely when
$[G:2G]<\infty$. We recall the two results behind this statement.

\begin{theorem}[{Charamaras--Mountakis~\cite[Corollary~1.13]{charamaras2024finding}}]
\label{G:thm:CM}
Let $(G,+)$ be a countable abelian group such that $2G$ is a subgroup of finite index, and
let $A\subset G$ have positive upper Banach density. Then there exist an infinite set
$B\subset A$ and some $t\in G$ such that $B\oplus B\subset A-t$.
\end{theorem}

(In~\cite{charamaras2024finding} the group is countable throughout, and, as the
conclusion shows, infinite. The result is cited as Corollary~1.12 in~\cite{kra2023problems}
and in~\cite[Theorem~1.1]{ackelsberg2024counterexamples}, following an earlier arXiv
version of~\cite{charamaras2024finding}.)

\begin{theorem}[{Ackelsberg~\cite[Theorem~1.2]{ackelsberg2024counterexamples}}]
\label{G:thm:Ack}
Suppose $\Gamma$ is a countable abelian group and $2\Gamma$ is a subgroup of infinite
index. Then for any $\eps>0$ there exists $A\subset\Gamma$ such that $\ubd(A)>1-\eps$
and such that if $B\subset\Gamma$ and $t\in\Gamma$ satisfy $B\oplus B+t\subset A$, then
$B$ is finite.
\end{theorem}

Ackelsberg also records~\cite[footnote~2]{ackelsberg2024counterexamples} that T.~Tao
observed independently that the condition $[\Gamma:2\Gamma]<\infty$ in \cref{G:thm:CM}
is necessary.

Conjecture~5.15 fails trivially for finite groups. Hence ``assuming that Conjecture~5.15
holds'' in \cref{G:q:516} means that $G$ is a countably infinite abelian group with
$[G:2G]<\infty$. We fix such a group for the rest of this section. For $A\subset G$ we let
\[
T(A)=\{t\in G:\ B\oplus B+t\subset A\ \text{for some infinite } B\subset G\}
\]
be the set of \emph{admissible translates} for $A$. We say that a set $T\subset G$
\emph{suffices} if $T(A)\cap T\neq\emptyset$ for every $A\subset G$ with $\ubd(A)>0$.

\subsection{Three elementary lemmas}

\begin{lemma}\label{G:lem:density}
Let $G$ be a countably infinite abelian group, and let $H\le G$ be a subgroup of finite
index $\ell$.
\begin{enumerate}[label=\textup{(\roman*)}]
\item For every \Fol{} sequence $\Phi$ and every $r\in G$ we have
$|(r+H)\cap\Phi_N|/|\Phi_N|\to1/\ell$. In particular $\ubd(r+H)=1/\ell>0$.
\item For all $X,Y\subset G$ we have $\ubd(X\cup Y)\le\ubd(X)+\ubd(Y)$.
\end{enumerate}
\end{lemma}

\begin{proof}
(i) Let $C\subset G$ and $x\in G$. Since $|(x+C)\cap\Phi_N|=|C\cap(\Phi_N-x)|$, we get
\[
\big||(x+C)\cap\Phi_N|-|C\cap\Phi_N|\big|\le|(\Phi_N-x)\,\triangle\,\Phi_N|
=2\big(|\Phi_N|-|\Phi_N\cap(\Phi_N-x)|\big)=o(|\Phi_N|).
\]
Let $r_1+H,\dots,r_\ell+H$ be the cosets of $H$. Applying this with $C=r_1+H$ and
$x=r_i-r_1$ shows that the $\ell$ quantities $|(r_i+H)\cap\Phi_N|$ differ pairwise by
$o(|\Phi_N|)$. They add up to $|\Phi_N|$, so each of them equals
$|\Phi_N|/\ell+o(|\Phi_N|)$.

(ii) For each \Fol{} sequence $\Phi$ we have
$\udens_\Phi(X\cup Y)\le\udens_\Phi(X)+\udens_\Phi(Y)\le\ubd(X)+\ubd(Y)$. Now take
the supremum over $\Phi$.
\end{proof}

\begin{lemma}[Pigeonhole]\label{G:lem:pigeon}
If $[G:2G]<\infty$ and $B\subset G$ is infinite, then $(B\oplus B)\cap2G\neq\emptyset$.
Consequently, if $B\oplus B+t\subset A$, then $t\in A+2G$. Moreover, if
$B\oplus B+t$ is contained in a single coset $r+2G$, then $t\in r+2G$.
\end{lemma}

\begin{proof}
Since $B$ is infinite and $2G$ has finitely many cosets, there are distinct $b,b'\in B$ in
the same coset $x+2G$, say $b=x+2u$ and $b'=x+2v$. Then $b+b'=2(x+u+v)\in2G$ and
$b+b'\in B\oplus B$. If $B\oplus B+t\subset A$, then $t+b+b'\in A$, so
$t\in A-(b+b')\subset A+2G$. If $B\oplus B+t\subset r+2G$, then $t+b+b'\in r+2G$, so
$t\in r+2G$.
\end{proof}

\begin{lemma}[Moving the translate]\label{G:lem:move}
For $B\subset G$ and $s\in G$ we have $(B-s)\oplus(B-s)=B\oplus B-2s$. Consequently
$T(A)+2G=T(A)$ for every $A\subset G$, i.e.\ $T(A)$ is a union of cosets of $2G$.
\end{lemma}

\begin{proof}
The map $b\mapsto b-s$ is a bijection $B\to B-s$, and $b-s\neq b'-s$ if and only if
$b\neq b'$; this gives the identity. If $B\oplus B+t\subset A$ with $B$ infinite, then
$B-s$ is infinite and $(B-s)\oplus(B-s)+(t+2s)=B\oplus B+t\subset A$. Thus
$t+2s\in T(A)$ for every $s\in G$.
\end{proof}

\subsection{The admissible translates}

\begin{theorem}\label{G:thm:translates}
Let $G$ be a countably infinite abelian group with $[G:2G]<\infty$. For $A\subset G$ put
$R(A)=\{r\in G:\ \ubd(A\cap(r+2G))>0\}$.
\begin{enumerate}[label=\textup{(\alph*)}]
\item\label{G:it:union} $T(A)$ is a union of cosets of $2G$.
\item\label{G:it:sandwich} $R(A)\subset T(A)\subset A+2G$. If $\ubd(A)>0$, then
$R(A)\neq\emptyset$, and hence $T(A)$ contains a full coset of $2G$.
\item\label{G:it:realise} For every nonempty union $U$ of cosets of $2G$ we have
$\ubd(U)>0$ and $T(U)=U$. Consequently, the family
$\{T(A):A\subset G,\ \ubd(A)>0\}$ is exactly the family of nonempty unions of cosets of
$2G$.
\item\label{G:it:suffice} A set $T\subset G$ suffices if and only if $T$ meets every coset of
$2G$. The minimal sets that suffice are exactly the transversals of $2G$ in $G$; they all
have exactly $[G:2G]$ elements.
\end{enumerate}
\end{theorem}

\begin{proof}
\ref{G:it:union} This is \cref{G:lem:move}.

\ref{G:it:sandwich} The inclusion $T(A)\subset A+2G$ is \cref{G:lem:pigeon}. Let
$r\in R(A)$ and put $A'=A\cap(r+2G)$. By \cref{G:thm:CM} applied to $A'$, there are an
infinite $B\subset G$ and $t\in G$ with $B\oplus B+t\subset A'\subset A$. Thus
$t\in T(A)$. Since $B\oplus B+t\subset r+2G$, \cref{G:lem:pigeon} gives $t\in r+2G$. By
part~\ref{G:it:union}, $r\in t+2G\subset T(A)$. Hence $R(A)\subset T(A)$.
If $\ubd(A)>0$, then since $A$ is the union of the finitely many sets $A\cap(r+2G)$,
\cref{G:lem:density}(ii) shows that $\ubd(A\cap(r+2G))>0$ for some $r$, i.e.\
$R(A)\neq\emptyset$. Since $R(A)$ is a union of cosets of $2G$, the last assertion follows.

\ref{G:it:realise} By \cref{G:lem:density}(i), $\ubd(U)\ge1/[G:2G]>0$. For every
$r\in U$ we have $U\cap(r+2G)=r+2G$, which has positive upper Banach density; hence
$U\subset R(U)$. By part~\ref{G:it:sandwich},
$U\subset R(U)\subset T(U)\subset U+2G=U$. The final assertion follows from this and
parts~\ref{G:it:union} and~\ref{G:it:sandwich}.

\ref{G:it:suffice} Suppose $T$ meets every coset of $2G$, and let $\ubd(A)>0$. By
part~\ref{G:it:sandwich}, $T(A)$ contains a coset of $2G$, and this coset meets $T$.
Conversely, suppose $T\cap(r+2G)=\emptyset$ for some $r$. Let $A=r+2G$. Then
$\ubd(A)>0$, and by part~\ref{G:it:realise}, $T(A)=r+2G$, which is disjoint from $T$.
So $T$ does not suffice.
A transversal meets every coset, so it suffices, and removing any of its points leaves
a coset unmet; hence transversals are minimal. If a sufficient set $T$ is not a
transversal, then it contains two points of one coset, and removing one of them leaves a
sufficient set. Hence $T$ is not minimal.
\end{proof}

\Cref{G:thm:translates}\ref{G:it:suffice} answers \cref{G:q:516}, and it confirms the
suspicion in \cref{G:q:53a} in the strongest possible form: the translates can be
restricted to any prescribed set of coset representatives of $2G$, and no smaller set
of translates works. Parts~\ref{G:it:union}--\ref{G:it:realise} concern an individual
set $A$. The admissible translates form a union of cosets of $2G$. This union contains
every coset in which $A$ has positive upper Banach density, and it is contained in the
union of the cosets that $A$ meets. Every nonempty union of cosets arises in this way.
In general neither bound is attained, as the next proposition shows. The only exact
statement for an individual $A$ is the tautological reduction to a single coset:
$r\in T(A)$ if and only if $A\cap(r+2G)\supset B\oplus B+t$ for some infinite $B$ and
some $t$. (If $B\oplus B+r\subset A$, pass to an infinite $B'\subset B$ inside one coset
of $2G$. Then $B'\oplus B'\subset 2G$, so $B'\oplus B'+r\subset A\cap(r+2G)$. The converse
follows from \cref{G:lem:pigeon,G:lem:move}.)

\begin{proposition}\label{G:prop:strict}
Let $G=\Z$.
\begin{enumerate}[label=\textup{(\roman*)}]
\item If $A=2\Z\cup\{3^n:n\in\N\}$, then $T(A)=2\Z$, while $A+2\Z=\Z$.
\item If $A=2\Z\cup X$ with $X=\{2^i+2^j+1:\ 1\le i<j\}$, then $R(A)=2\Z$, while
$T(A)=\Z$.
\end{enumerate}
So both inclusions in \cref{G:thm:translates}\ref{G:it:sandwich} can be strict.
\end{proposition}

\begin{proof}
(i) By \cref{G:thm:translates}\ref{G:it:sandwich}, $2\Z\subset R(A)\subset T(A)$. Since
$3\in A$, we have $A+2\Z=\Z$. Suppose that $t$ is odd and $B\oplus B+t\subset A$ with $B$
infinite. Choose an infinite $B'\subset B$ whose elements all have the same parity, and fix
$b_1<b_2$ in $B'$. For every $b\in B'\setminus\{b_1,b_2\}$ the numbers $b_1+b+t$ and
$b_2+b+t$ are odd elements of $A$, hence powers of $3$, and they differ by $b_2-b_1$. The
equation $3^m-3^n=b_2-b_1$ with $m>n\ge1$ has only finitely many solutions, because
$3^n\le b_2-b_1$ and $n$ determines $m$. But $b_1+b+t$ takes infinitely many values as $b$
ranges over $B'$. This is a contradiction, so $T(A)=2\Z$.

(ii) For $B=\{2^i:i\in\N\}$ we have $B\oplus B+1=X\subset A$, so $1\in T(A)$, and
$T(A)=\Z$ by \cref{G:lem:move}. Clearly $2\Z\subset R(A)$. It remains to show that
$\ubd(X)=0$. We claim that if every interval of $L$ consecutive integers contains at most
$c$ elements of a set $Y\subset\Z$, then $\ubd(Y)\le c/L$. Indeed, for a \Fol{} sequence
$\Phi$ we have
\[
\sum_{x=0}^{L-1}|Y\cap(\Phi_N+x)|=\sum_{y\in\Phi_N}|Y\cap\{y,\dots,y+L-1\}|\le c|\Phi_N|,
\]
and $|Y\cap(\Phi_N+x)|\ge|Y\cap\Phi_N|-|\Phi_N\setminus(\Phi_N+x)|=|Y\cap\Phi_N|-o(|\Phi_N|)$.
Hence $L\,\udens_\Phi(Y)\le c$.

Now let $I$ be an interval of $L\ge2$ consecutive integers. An element $2^i+2^j+1$ of
$X$ with $i<j$ lies in $(2^j,2^{j+1})$. The ranges $[2^j,2^{j+1})$ with $2^j\ge L$ have
length at least $L$, so at most two of them meet $I$. There are at most $\log_2L+1$ ranges
with $2^j<L$. For fixed $j$, the elements of $X\cap I$ correspond to distinct powers $2^i$
in the interval $I-2^j-1$ of length $L$. Such an interval contains at most one power of $2$
that is $\ge L$, and at most $\log_2L+1$ powers of $2$ that are $<L$. Hence
$|X\cap I|\le(\log_2L+3)^2$, and so $\ubd(X)\le(\log_2L+3)^2/L$ for every $L\ge2$. Thus
$\ubd(X)=0$ and $R(A)=2\Z$.
\end{proof}

\formalization{All four parts of \cref{G:thm:translates} are formalized in the file
\leanfile{Q516} as \leanname{thm_translates}, with \cref{G:thm:CM} stated as an explicit
hypothesis. The statements that do not use it, and \cref{G:prop:strict}
(\leanname{prop_strict_i}, \leanname{prop_strict_ii}), are proved unconditionally.}

\begin{remark}\label{G:rem:examples}
We compare \cref{G:thm:translates} with the examples in~\cite[\S5.3]{kra2023problems}.
\begin{enumerate}[label=\textup{(\arabic*)}]
\item For $G=\Z$ we have $2G=2\Z$, and the transversal $\{0,1\}$ suffices. This matches
the case $k=2$ of~\cite[Proposition~2.6]{kra2023problems}, which is stated in $\N$.
Taking $r=0$ in part~\ref{G:it:sandwich}, we get that $t=0$ is admissible whenever
$\ubd(A\cap2G)>0$. This is the group analogue of item~3
of~\cite[Proposition~2.6]{kra2023problems} for $k=2$. That proposition is an equivalence
between statements for all $k$, but its proof works for each $k$ separately. So for $k=2$
its items hold unconditionally, by~\cite[Theorem~2.2]{kra2023problems}.
\item If $G=\bigoplus_{i\in\N}\Z/3\Z$, then $2G=G$, so $T=\{0\}$ suffices, as observed
in~\cite{kra2023problems}.
\item If $G=\bigoplus_{i\in\N}\Z/2\Z$, then $2G=\{0\}$ has infinite index. This group lies
outside the scope of \cref{G:q:516}, because by \cref{G:thm:Ack} Conjecture~5.15 fails
there. In groups with $1<[G:2G]<\infty$, the phenomenon pointed out
in~\cite{kra2023problems} for this group (one cannot always take $t=0$) is explained by
cosets: for $A=r+2G$ with $r\notin2G$, every admissible translate lies in $r+2G$.
\end{enumerate}
\end{remark}

\begin{remark}[Credit]\label{G:rem:credit}
The only deep ingredient is \cref{G:thm:CM}. The two elementary devices in
\cref{G:lem:pigeon,G:lem:move} are not new. For unrestricted sumsets $t+B+B$,
Charamaras, Kousek, Mountakis and Radi\'c~\cite[Lemmas~2.1 and~2.2]{charamaras2025infinite}
use pigeonholing into cosets of $2G$ and write $t=2s+g_i$ with $g_i$ a coset
representative, replacing $B$ by $B+s$. Ackelsberg and Jamneshan~\cite[discussion after
Theorem~1.9, and Corollary~1.10]{ackelsberg2025equidistribution} use the same two
arguments for $B\oplus B\oplus B$ and the cosets of $3\Gamma$, under the assumption
$[\Gamma:6\Gamma]<\infty$. For $G=\Z$ the reduction
appears in~\cite[Proposition~2.6]{kra2023problems}. Our contribution consists of the exact
characterisation of the sets of translates that suffice, the minimality of transversals,
the description of the possible sets $T(A)$, and the two-sided bounds for $T(A)$ for
$B\oplus B$.
\end{remark}

% ---------------------------------------------------------------------------

\section{Cartesian products in \texorpdfstring{$\N^2$}{N2}}
\label{G:sec:q510}

In this section we answer Question~5.10 of~\cite{kra2023problems} negatively. We
construct sets $A\subset\N^2$ all of whose fibres have positive density along a fixed
\Fol{} sequence (in one version, every fibre has natural density exactly $1/4$), such that
no translate of $A$ contains the Ramsey configuration
$\{(b_1,b_2):b_1,b_2\in B,\ b_1<b_2\}$ for an infinite $B\subset\N$.

\subsection{The question}
Throughout, $\N=\{1,2,\dots\}$. Following~\cite{kra2023problems}, a \emph{\Fol{} sequence
in $\N$} is a sequence $\Phi=(\Phi_N)$ of finite subsets of $\N$ with
$|(\Phi_N+t)\cap\Phi_N|/|\Phi_N|\to1$ for every $t\in\N$. The upper and lower densities
along $\Phi$ are $\udens_\Phi$ and $\ldens_\Phi$. When they are equal, their common value
is denoted $\dens_\Phi$. For $A\subset\N^2$ and $n\in\N$ the \emph{fibre} of $A$ over $n$ is
$A_n=\{m\in\N:(n,m)\in A\}$. The configuration in question is
\begin{equation}\label{G:eq:BBcart}
\{(b_1,b_2):\ b_1,b_2\in B,\ b_1<b_2\}\subset A-t,
\end{equation}
which is equation~(21) of~\cite{kra2023problems}. It says that
$(b_1+t_1,b_2+t_2)\in A$ for all $b_1<b_2$ in $B$, where $t=(t_1,t_2)$.

\begin{question}{5.10}\label{G:q:510}
Suppose $A\subset\N^2$ satisfies
\begin{equation}\label{G:eq:22}
\liminf_{n\in\N}\dens_\Phi(A_n)>0
\end{equation}
for some \Fol{} sequence $\Phi$ on $\N$. Does there exist an infinite set $B\subset\N$ and
an element $t\in\N^2$ satisfying~\eqref{G:eq:BBcart}?
\end{question}

The survey points out that \eqref{G:eq:BBcart} can only hold if infinitely many fibres
$A_n$ have the property that every finite subfamily has infinite intersection. Example~5.11
of~\cite{kra2023problems}, due to F.~Hern\'andez Castro, takes $A=X\times Y$ with
$X=\bigcup_n[4^n,c4^n)$ and $Y=\bigcup_n[c'4^n,c''4^n)$ (intersected with $\N$), where
$1<c<c'<c''<2$. Then $X\cap(Y-t)$ is finite for every $t\ge0$, and hence no translate of
$A$ contains a configuration~\eqref{G:eq:BBcart}. In that example the fibre $A_n$ is $Y$
for $n\in X$ and empty for $n\notin X$. So it shows that the $\liminf$ in~\eqref{G:eq:22}
cannot be replaced by a $\limsup$. It does not satisfy~\eqref{G:eq:22}. We modify the
example so that \emph{every} fibre has positive density. Two products of this type are
glued along a partition of the rows, and the second coordinate of each product is
offset from its own row set.

\subsection{A criterion}

\begin{lemma}\label{G:lem:criterion}
Let $\N=P_0\sqcup P_1$ be a partition, and let $S_0,S_1\subset\N$ be sets such that
\begin{equation}\label{G:eq:finite}
S_c\cap(P_c+s)\ \text{is finite for every } s\in\Z\text{ and } c\in\{0,1\}.
\end{equation}
Let $A=(P_0\times S_0)\cup(P_1\times S_1)$. Then $A_n=S_c$ for every $n\in P_c$. Moreover,
there are no infinite set $B\subset\N$ and $t=(t_1,t_2)\in\Z^2$ such that
$(b_1+t_1,b_2+t_2)\in A$ for all $b_1<b_2$ in $B$. In particular, $A$ satisfies
\eqref{G:eq:BBcart} for no infinite $B\subset\N$ and no $t\in\N^2$.
\end{lemma}

\begin{proof}
The identity $A_n=S_c$ for $n\in P_c$ holds because $P_0\cap P_1=\emptyset$. Suppose that
$B$ and $t$ exist. Every $b\in B$ is the smaller element of some pair in $B$, so
$b+t_1\in\N$. By the pigeonhole principle there are $c\in\{0,1\}$ and an infinite
$B'\subset B$ with $b+t_1\in P_c$ for all $b\in B'$. Let $b_0=\min B'$. For every
$b\in B'\setminus\{b_0\}$ we have $(b_0+t_1,b+t_2)\in A$ and $b_0+t_1\in P_c$, so
$b+t_2\in A_{b_0+t_1}=S_c$. Also $b+t_1\in P_c$, so
$b+t_2\in P_c+(t_2-t_1)$. Hence the infinitely many distinct numbers $b+t_2$, for
$b\in B'\setminus\{b_0\}$, lie in $S_c\cap(P_c+(t_2-t_1))$. This contradicts
\eqref{G:eq:finite}.
\end{proof}

\subsection{Two examples}

\begin{example}[Scale $4^k$]\label{G:ex:4k}
Let
\[
P_0=\bigcup_{k\ge0}[4^k,2\cdot4^k),\qquad P_1=\bigcup_{k\ge0}[2\cdot4^k,4^{k+1}),
\]
\[
S_0=\bigcup_{k\ge1}[9\cdot4^{k-1},15\cdot4^{k-1}),\qquad
S_1=\bigcup_{k\ge1}[5\cdot4^{k-1},7\cdot4^{k-1}),
\]
all intervals being intervals of integers. So $\N=P_0\sqcup P_1$. Each block of $S_0$
lies in the middle of the block $[2\cdot4^k,4^{k+1})$ of $P_1$, and each block of $S_1$
lies in the middle of the block $[4^k,2\cdot4^k)$ of $P_0$. Put
$A=(P_0\times S_0)\cup(P_1\times S_1)$.
\end{example}

\begin{example}[Square blocks]\label{G:ex:blocks}
For $k\in\N$ let $J_k=[k^2,(k+1)^2)\cap\N$, so that $\N=\bigsqcup_{k\ge1}J_k$ and
$|J_k|=2k+1$. Let $P'_c=\bigcup_{k\equiv c\ (\mathrm{mod}\ 2)}J_k$ for $c\in\{0,1\}$, and let
$S'_c$ consist of the \emph{middle parts} of the blocks of the opposite parity:
\[
S'_c=\{k^2+r:\ k\in\N,\ k\equiv1-c\ (\mathrm{mod}\ 2),\ r\in\Z,\ k\le2r<3k\}.
\]
Since $k\le 2r<3k$ forces $0\le r\le 2k$, we have $S'_c\cap J_k=\{k^2+r:k\le2r<3k\}$ when
$k\equiv1-c$, and this set has exactly $k$ elements. Put
$A'=(P'_0\times S'_0)\cup(P'_1\times S'_1)$.
\end{example}

\begin{proposition}\label{G:prop:examples}
\begin{enumerate}[label=\textup{(\roman*)}]
\item The sets of \cref{G:ex:4k} satisfy \eqref{G:eq:finite}. More precisely, every element
of $S_c\cap(P_c+s)$ is smaller than $15|s|$. Along the \Fol{} sequence
$\Phi_N=[1,4^N]$ we have $\dens_\Phi(S_0)=1/2$ and $\dens_\Phi(S_1)=1/6$. Hence every fibre
of $A$ has density $1/2$ or $1/6$ along $\Phi$, and $A$ satisfies~\eqref{G:eq:22}.
\item The sets of \cref{G:ex:blocks} satisfy \eqref{G:eq:finite}. More precisely, every
element of $S'_c\cap(P'_c+s)$ is smaller than $(2|s|+1)^2$. Both $S'_0$ and $S'_1$ have
natural density $1/4$. Hence every fibre of $A'$ has natural density $1/4$, and $A'$
satisfies \eqref{G:eq:22} along $\Phi_N=[1,N]$.
\end{enumerate}
\end{proposition}

\begin{proof}
(i) Let $m\in S_0$, say $9\cdot4^{k-1}\le m<15\cdot4^{k-1}$, and suppose $|s|<4^{k-1}$.
Then $2\cdot4^k=8\cdot4^{k-1}<m-s<16\cdot4^{k-1}=4^{k+1}$, so $m-s\in P_1$ and
$m-s\notin P_0$. Similarly, if $m\in S_1$, say $5\cdot4^{k-1}\le m<7\cdot4^{k-1}$, and
$|s|<4^{k-1}$, then $4^k<m-s<2\cdot4^k$, so $m-s\in P_0$ and $m-s\notin P_1$. Hence
every $m\in S_c\cap(P_c+s)$ lies in a block with $4^{k-1}\le|s|$, and so
$m<15\cdot4^{k-1}\le15|s|$. (For $s=0$ the intersection is empty.)

The sets $[1,4^N]$ form a \Fol{} sequence in $\N$, since $|([1,M]+t)\cap[1,M]|=M-t$ for
$0\le t\le M$. For $N\ge2$ the blocks of $S_0$ and $S_1$ with $k\le N-1$ lie in
$[1,4^N]$, since $15\cdot4^{N-2}<4^N$. The blocks with $k\ge N$ start at
$9\cdot4^{N-1}$ and $5\cdot4^{N-1}$, both larger than $4^N$, so they are disjoint from
$[1,4^N]$. Therefore
\[
|S_0\cap[1,4^N]|=\sum_{k=1}^{N-1}6\cdot4^{k-1}=2(4^{N-1}-1),\qquad
|S_1\cap[1,4^N]|=\sum_{k=1}^{N-1}2\cdot4^{k-1}=\tfrac23(4^{N-1}-1),
\]
and dividing by $4^N$ gives the limits $1/2$ and $1/6$. Every fibre $A_n$ equals $S_0$ or
$S_1$ by \cref{G:lem:criterion}, so $\inf_n\dens_\Phi(A_n)=1/6>0$.

(ii) Let $m=k^2+r\in S'_c$ with $k\le2r<3k$ and $k\equiv1-c\pmod2$. Suppose $|s|<k/2$.
Then $r-s\ge r-|s|>r-k/2\ge0$ and $r-s\le r+|s|<3k/2+k/2=2k$. Hence
$m-s=k^2+(r-s)\in J_k\subset P'_{1-c}$, so $m-s\notin P'_c$. Consequently every element
of $S'_c\cap(P'_c+s)$ lies in a block $J_k$ with $k\le2|s|$, and hence is smaller than
$(2|s|+1)^2$.

For the density, let $N\ge1$ and let $K\ge1$ be such that $K^2\le N<(K+1)^2$. The blocks
$J_k$ with $k<K$ lie in $[1,N]$, and $J_K$ is the only block that meets $[1,N]$ without
lying inside it. Hence
\[
|S'_c\cap[1,N]|=\sum_{\substack{1\le k<K\\ k\equiv1-c\ (2)}}k+O(K)
=\frac{K^2}{4}+O(K)=\frac N4+O(\sqrt N).
\]
So $\dens(S'_c)=1/4$. Every fibre of $A'$ is $S'_0$ or $S'_1$, and $[1,N]$ is a \Fol{}
sequence.
\end{proof}

\begin{theorem}\label{G:thm:q510}
The sets $A$ of \cref{G:ex:4k} and $A'$ of \cref{G:ex:blocks} satisfy~\eqref{G:eq:22}. In
fact every fibre has a density along the indicated \Fol{} sequence: $1/2$ or $1/6$
for $A$, and natural density exactly $1/4$ for $A'$. Nevertheless, there is no infinite
$B\subset\N$ and no $t\in\Z^2$ such that $(b_1+t_1,b_2+t_2)$ lies in $A$ (respectively
$A'$) for all $b_1<b_2$ in $B$. In particular, the answer to \cref{G:q:510} is negative.
\end{theorem}

\begin{proof}
Combine \cref{G:prop:examples} with \cref{G:lem:criterion}.
\end{proof}

\formalization{\Cref{G:lem:criterion}, \cref{G:prop:examples} and \cref{G:thm:q510} are
formalized in \leanfile{Q510} (\leanname{criterion}, \leanname{q510_example_4k},
\leanname{q510_example_blocks}, \leanname{q510_negative}).}

\begin{remark}\label{G:rem:q510}
\begin{enumerate}[label=\textup{(\arabic*)}]
\item The set $A'$ shows that the answer is negative already for the standard \Fol{}
sequence $[1,N]$, that is, with natural densities: every fibre has natural density $1/4$.
It does not satisfy~\eqref{G:eq:22} along every \Fol{} sequence. Each $S'_c$ misses the
blocks $J_k$ with $k\equiv c$, and these have length $2k+1\to\infty$. So along
$\Phi_N=J_{2N}$, for instance, the fibres $S'_0$ have density $0$. The hypothesis that
\eqref{G:eq:22} holds along every \Fol{} sequence, i.e.\ that the fibres have uniformly
positive lower Banach density, is not addressed here; compare~(4). The set $A$ of \cref{G:ex:4k} has fibres without natural density. Its fibres have
lower natural densities $2/9$ and $2/15$, which are approached at the integers just
below the left endpoints of the blocks of $S_0$ and $S_1$, respectively.
\item The survey explains that a positive answer to \cref{G:q:510} would imply
Theorem~2.2 and Conjecture~3.10 of~\cite{kra2023problems}. A negative answer has no
bearing on either statement.
\item The survey's Theorem~5.12 shows that a weaker hypothesis gives infinite $B,C$ with
$\{(b,c):b\in B,\ c\in C,\ b<c\}\subset A$. This is consistent with
\cref{G:thm:q510}: our sets contain even the full products $P_0\times S_0$ and
$P'_0\times S'_0$. What fails is the diagonal coupling in~\eqref{G:eq:BBcart}: every
element of $B$, except the smallest, must serve both as a row index and as a column index.
\item The construction in \cref{G:lem:criterion} cannot work if every fibre is required
to be syndetic with a gap bound that does not depend on the fibre. Indeed, suppose that for each $c$ every
$n\in\N$ satisfies $n+s\in S_c$ for some $0\le s\le L$. Then
$P_c\subset\bigcup_{0\le s\le L}(S_c-s)$, and \eqref{G:eq:finite} fails for the class $c$
with $P_c$ infinite. We do not know the answer to \cref{G:q:510} under this stronger
hypothesis.
\end{enumerate}
\end{remark}

% ---------------------------------------------------------------------------

\section{Corners in \texorpdfstring{$G^2$}{G2} and \texorpdfstring{$G^d$}{Gd}}
\label{G:sec:q523}

In this section we answer Question~5.23 of~\cite{kra2023problems} and its $G^d$ version,
Question~\S5.4(b), negatively. The negative answer to the question as stated is elementary.
The hypothesis controls only the fibres in the last coordinate, while the conclusion
requires infinitely many points of $A$ on a single line in the first coordinate
direction. So this part of the answer depends on the formulation. We therefore also
treat the natural repaired version in which these points are dropped. Using
\cref{G:sec:q510}, we show that it fails as well for $G=\Z$ (and for $G=\N$). The
repair that adds a positive-density hypothesis on the columns remains open.

\subsection{The question}
Let $(G,\cdot)$ be a group. Following~\cite[\S5.4]{kra2023problems}, a \emph{left \Fol{}
sequence} is a sequence $(\Phi_N)$ of finite nonempty subsets of $G$ with
$|\Phi_N\cap g\Phi_N|/|\Phi_N|\to1$ for every $g\in G$. Austin's theorem on corners in
amenable groups~\cite{austin2016nonconventional}, quoted
in~\cite[Theorem~5.22]{kra2023problems}, motivates the following question.

\begin{question}{5.23}\label{G:q:523}
Let $(G,\cdot)$ be an amenable group and let $N\mapsto\Phi_N$ be a left \Fol{} sequence in
$G$. If $A\subset G^2$ satisfies
\begin{equation}\label{G:eq:23}
\liminf_{g\to\infty}\ \liminf_{M\to\infty}\ \frac1{|\Phi_M|}\sum_{h\in\Phi_M}\one_A(g,h)>0,
\end{equation}
where $g\to\infty$ means leaving all finite sets, do there exist $(t,s)\in G^2$ and an
injective sequence $b\colon\N\to G$ such that
\begin{equation}\label{G:eq:corner}
\{(t,s),\ (b(i)t,s),\ (b(i)t,b(j)s):\ i<j\}\subset A\,?
\end{equation}
\end{question}

The survey adds that if $G=\N$, an affirmative answer to \cref{G:q:523} implies an
affirmative answer to \cref{G:q:510}. Two remarks on the setting are in order. First, a
group with a left \Fol{} sequence is countable. Indeed, let $H$ be the subgroup generated by
$\bigcup_N\Phi_N$, which is countable. If $g\notin H$, then $g\Phi_N\subset gH$ is disjoint
from $\Phi_N\subset H$, so the \Fol{} condition fails for $g$. Second, if $G$ is finite,
there is no injective $b\colon\N\to G$. So we assume, as the survey intends, that $G$ is
countably infinite.

For $A\subset G^2$ and $g\in G$ write $A_g=\{h:(g,h)\in A\}$ for the \emph{row fibre}
and $A^h=\{g:(g,h)\in A\}$ for the \emph{column}. Condition~\eqref{G:eq:23} concerns only
the row fibres. The points $(b(i)t,s)$, $i\in\N$, in \eqref{G:eq:corner} are pairwise
distinct (since $b$ is injective) and lie in the single column $A^s$.

\begin{lemma}\label{G:lem:folnergrow}
If $G$ is an infinite group and $(\Phi_M)$ is a left \Fol{} sequence, then
$|\Phi_M|\to\infty$.
\end{lemma}

\begin{proof}
Suppose that $|\Phi_M|\le K$ for infinitely many $M$. Choose a finite $S\subset G$ with
$|S|>K$. For each $x\in S$ the number $|\Phi_M\setminus x\Phi_M|$ is a nonnegative
integer, and $|\Phi_M\setminus x\Phi_M|/|\Phi_M|\to0$. Along the $M$ with $|\Phi_M|\le K$
we therefore have $|\Phi_M\setminus x\Phi_M|=0$ for all $x\in S$ once $M$ is large. Then
$\Phi_M\subset x\Phi_M$, and since $|x\Phi_M|=|\Phi_M|$, we get $x\Phi_M=\Phi_M$. Pick
such an $M$ and some $y\in\Phi_M$. Then $Sy\subset\Phi_M$, so $|\Phi_M|\ge|S|>K$, a
contradiction.
\end{proof}

\subsection{The question as stated, in \texorpdfstring{$G^2$}{G2} and
\texorpdfstring{$G^d$}{Gd}}

The survey states that \cref{G:q:523} ``can be generalized to $G^d$ for $d>2$ in a
straightforward manner'', but does not write out the generalisation. Modelled on
Austin's $d$-dimensional corners~\cite[Theorem~5.22]{kra2023problems}, the natural
conclusion asks for $t\in G^d$ and an injective $b\colon\N\to G$ such that, for all
$i_1<\dots<i_d$, the $d+1$ points
\begin{equation}\label{G:eq:cornerd}
(t_1,\dots,t_d),\ (b(i_1)t_1,t_2,\dots,t_d),\ (b(i_1)t_1,b(i_2)t_2,t_3,\dots,t_d),\ \dots,\
(b(i_1)t_1,\dots,b(i_d)t_d)
\end{equation}
lie in $A$. The natural analogues of~\eqref{G:eq:23} constrain the fibres in the last
coordinate,
\begin{equation}\label{G:eq:I}
\liminf_{(x_1,\dots,x_{d-1})\to\infty}\ \liminf_{M\to\infty}\
\frac{|\{h\in\Phi_M:(x_1,\dots,x_{d-1},h)\in A\}|}{|\Phi_M|}>0,
\end{equation}
or the row fibres over the first coordinate,
\begin{equation}\label{G:eq:II}
\liminf_{x_1\to\infty}\ \liminf_{M\to\infty}\
\frac{|\{y\in\Phi_M^{d-1}:(x_1,y)\in A\}|}{|\Phi_M|^{d-1}}>0.
\end{equation}

\begin{question}{\S5.4(b)}\label{G:q:54b}
Let $d\ge3$, let $G$ be a countably infinite amenable group with a left \Fol{} sequence
$\Phi$, and suppose $A\subset G^d$ satisfies \eqref{G:eq:I} or \eqref{G:eq:II}. Do there
exist $t\in G^d$ and an injective $b\colon\N\to G$ such that the points
\eqref{G:eq:cornerd} lie in $A$ for all $i_1<\dots<i_d$?
\end{question}

For $d=2$ both \eqref{G:eq:I} and \eqref{G:eq:II} reduce to \eqref{G:eq:23}, and
\cref{G:q:54b} reduces to \cref{G:q:523}.

\begin{theorem}\label{G:thm:q523literal}
Let $G$ be a countably infinite group, enumerated as $G=\{g_1,g_2,\dots\}$ without
repetitions, and let $d\ge2$. Put
\[
E_d=\{x\in G^d:\ x_1=g_n,\ x_d=g_m\ \text{for some } n<m\}.
\]
Let $\Phi$ be any left \Fol{} sequence in $G$. Then:
\begin{enumerate}[label=\textup{(\roman*)}]
\item for every $(x_1,\dots,x_{d-1})\in G^{d-1}$, the last-coordinate fibre
$\{h:(x_1,\dots,x_{d-1},h)\in E_d\}$ is cofinite, and its density along $\Phi$ equals $1$;
\item for every $x_1\in G$, the set $\{y\in G^{d-1}:(x_1,y)\in E_d\}$ has density $1$
along $\Phi^{d-1}=(\Phi_M^{d-1})_M$;
\item $|E_d\cap\Phi_M^d|/|\Phi_M|^d\to1/2$;
\item for every $t\in G^d$, the set $\{g\in G:(gt_1,t_2,\dots,t_d)\in E_d\}$ is finite.
\end{enumerate}
In particular $E_d$ satisfies~\eqref{G:eq:I} and~\eqref{G:eq:II}, and
\eqref{G:eq:23} when $d=2$, with all limits equal to $1$. But there is no $t\in G^d$
and no injective $b$ such that $(b(i)t_1,t_2,\dots,t_d)\in E_d$ for infinitely many $i$.
Hence the answers to \cref{G:q:523} and \cref{G:q:54b} are negative, for every countably
infinite amenable group $G$ and every left \Fol{} sequence.
\end{theorem}

\begin{proof}
If $x_1=g_n$, then the last-coordinate fibre is $G\setminus\{g_1,\dots,g_n\}$, whatever
$x_2,\dots,x_{d-1}$ are. Its density along $\Phi$ is at least $1-n/|\Phi_M|$, which tends
to $1$ by \cref{G:lem:folnergrow}. This proves (i). Part (ii) follows similarly: the set in
question is $G^{d-2}\times(G\setminus\{g_1,\dots,g_n\})$, and its relative density in
$\Phi_M^{d-1}$ is at least $1-n/|\Phi_M|$. For (iii), note that among the
$|\Phi_M|^2-|\Phi_M|$ pairs $(x,y)\in\Phi_M^2$ with $x\neq y$, exactly one of $(x,y)$ and
$(y,x)$ has the first entry earlier in the enumeration. Hence
$|E_d\cap\Phi_M^d|=\frac12(|\Phi_M|^2-|\Phi_M|)\,|\Phi_M|^{d-2}$, and (iii) follows from
\cref{G:lem:folnergrow}. For (iv), let $t_d=g_m$. If $(gt_1,t_2,\dots,t_d)\in E_d$, then
$gt_1\in\{g_1,\dots,g_{m-1}\}$, so $g$ lies in the finite set
$\{g_1,\dots,g_{m-1}\}t_1^{-1}$. Since $b$ is injective, $(b(i)t_1,t_2,\dots,t_d)\in E_d$
can hold for at most $m-1$ values of $i$.
\end{proof}

\formalization{\Cref{G:thm:q523literal} and the resulting negative answers to
\cref{G:q:523,G:q:54b} are formalized in \leanfile{Q523} (\leanname{q523literal},
\leanname{q523_negative}). \Cref{G:thm:repairA} below is not formalized.}

For $G=\Z$ with $\Phi_M=[1,M]$, the simpler set $\{(g,h)\in\Z^2:h>|g|\}$ also works. For
$G=\N$ one can take $\{(n,m)\in\N^2:n<m\}$. In the semigroup $\N$, where the survey's
implication to \cref{G:q:510} is stated, this means that the answer to the literal
question is negative for a trivial reason.

\begin{remark}[The answer depends on the formulation]\label{G:rem:formulation}
\Cref{G:thm:q523literal} has nothing to do with sumsets: the conclusion
\eqref{G:eq:corner} contains infinitely many distinct points of a single column, and
\eqref{G:eq:23} does not even force columns to be infinite. The survey's comment that, for
$G=\N$, \cref{G:q:523} implies \cref{G:q:510} shows that the authors are mainly
interested in the points $(b(i)t,b(j)s)$. The points $(t,s)$ and $(b(i)t,s)$ come from
the shape of Austin's corners. There are two natural repairs:
\begin{enumerate}[label=\textup{(\alph*)}]
\item drop the points $(t,s)$ and $(b(i)t,s)$, and ask only for
$\{(b(i)t,b(j)s):i<j\}\subset A$ (and, in $G^d$, only for the top points
$(b(i_1)t_1,\dots,b(i_d)t_d)$ with $i_1<\dots<i_d$);
\item keep \eqref{G:eq:corner} but add to \eqref{G:eq:23} the analogous condition on
the columns, that is,
\[
\liminf_{h\to\infty}\ \liminf_{M\to\infty}\ \frac{|A^h\cap\Phi_M|}{|\Phi_M|}>0 .
\]
\end{enumerate}
In \cref{G:thm:repairA} below we show, using \cref{G:sec:q510}, that repair~(a) still has
a negative answer for $G=\Z$. Repair~(b), with positive-density hypotheses on both the
row fibres and the columns, is to our knowledge open. Neither the
examples of \cref{G:thm:q523literal} nor those of \cref{G:thm:repairA} satisfy it: for
instance, the columns of $\tilde A$ below are empty for $|h|\notin S'_0\cup S'_1$.
\end{remark}

\subsection{The repaired question (a)}

For $G=\N$, repair~(a) of \cref{G:q:523} is answered by \cref{G:thm:q510} directly.
Suppose that $t,s\in\N$ and an injective $b\colon\N\to\N$ satisfy
$(b(i)+t,b(j)+s)\in A'$ for all $i<j$. Since $b$ is injective, it is unbounded, so it has
a strictly increasing subsequence. Its range $B$ then satisfies \eqref{G:eq:BBcart} with
the translate $(t,s)$. This is impossible. Also, $A'$ satisfies \eqref{G:eq:23} with
$\Phi_M=[1,M]$, because every fibre of $A'$ has natural density $1/4$. To treat the group
$\Z$ and the $d$-dimensional version, we transfer $A'$ to $\Z^2$. For $x\in\Z$ let
$|x|_+=\max\{|x|,1\}$, and let
\[
\tilde A=\{(g,h)\in\Z^2:\ h\neq0,\ (|g|_+,|h|)\in A'\},
\]
where $A'=(P'_0\times S'_0)\cup(P'_1\times S'_1)$ is the set of \cref{G:ex:blocks}.

\begin{lemma}\label{G:lem:Ztransfer}
For every $g\in\Z$, the row fibre $\{h:(g,h)\in \tilde A\}$ meets $[1,M]$ in
$|S'_c\cap[1,M]|$ points, where $c$ is the class with $|g|_+\in P'_c$. Hence it has
density $1/4$ along $\Phi_M=[1,M]$. There are no $(t,s)\in\Z^2$ and injective
$b\colon\N\to\Z$ with $(b(i)+t,b(j)+s)\in \tilde A$ for all $i<j$.
\end{lemma}

\begin{proof}
The fibre over $g$ is $\{h\neq0:|h|\in S'_c\}$, and its intersection with $[1,M]$ is
$S'_c\cap[1,M]$. The density statement now follows from
\cref{G:prop:examples}(ii).

Suppose that $t,s,b$ exist. Since $b$ is injective, for every $L$ only finitely many $i$
satisfy $|b(i)|\le L$. Choose $\epsilon\in\{1,-1\}$ such that $\epsilon b(i)>0$ for
infinitely many $i$. Then choose $i_1<i_2<\cdots$ inductively so that the numbers
$\beta_k=\epsilon b(i_k)$ are positive and strictly increasing, with
$\beta_1>|t|+|s|+1$. For $k<l$ we have $(b(i_k)+t,b(i_l)+s)\in \tilde A$. Now
$|b(i_k)+t|_+=|\epsilon\beta_k+t|=\beta_k+\epsilon t$, because $\beta_k>|t|+1$, and similarly
$|b(i_l)+s|=\beta_l+\epsilon s$. So $(\beta_k+\epsilon t,\beta_l+\epsilon s)\in A'$ for all
$k<l$. Then the infinite set $B=\{\beta_k\}\subset\N$ and the translate
$(\epsilon t,\epsilon s)\in\Z^2$ contradict \cref{G:thm:q510}.
\end{proof}

\begin{theorem}\label{G:thm:repairA}
Let $d\ge2$ and $\tilde A^{(d)}=\{x\in\Z^d:(x_1,x_d)\in \tilde A\}$, and let $\Phi_M=[1,M]$.
\begin{enumerate}[label=\textup{(\roman*)}]
\item For every $(x_1,\dots,x_{d-1})\in\Z^{d-1}$, the last-coordinate fibre of $\tilde A^{(d)}$ has
density $1/4$ along $\Phi$. For every $x_1\in\Z$, the fibre
$\{y\in\Z^{d-1}:(x_1,y)\in \tilde A^{(d)}\}$ has density $1/4$ along $\Phi^{d-1}$. In particular
$\tilde A^{(d)}$ satisfies \eqref{G:eq:I} and \eqref{G:eq:II}, and $\tilde A^{(2)}=\tilde A$ satisfies
\eqref{G:eq:23}.
\item There are no $t\in\Z^d$ and injective $b\colon\N\to\Z$ such that
$(b(i_1)+t_1,\dots,b(i_d)+t_d)\in \tilde A^{(d)}$ for all $i_1<\dots<i_d$.
\end{enumerate}
Consequently repair~(a) of \cref{G:q:523} has a negative answer for $G=\Z$, and so does
its analogue for \cref{G:q:54b}, in which only the top points of \eqref{G:eq:cornerd}
are required.
\end{theorem}

\begin{proof}
(i) The last-coordinate fibre over $(x_1,\dots,x_{d-1})$ is the row fibre of $\tilde A$ over
$x_1$. The fibre over $x_1$ is $\Z^{d-2}\times\{h:(x_1,h)\in \tilde A\}$, and its relative
density in $[1,M]^{d-1}$ equals that of the row fibre in $[1,M]$. Now apply
\cref{G:lem:Ztransfer}.

(ii) Suppose that $t$ and $b$ exist. If $i<j$ and $j-i\ge d-1$, there are indices
$i<i_2<\dots<i_{d-1}<j$, and so $(b(i)+t_1,b(j)+t_d)\in \tilde A$. Let
$b'(k)=b((d-1)k)$. Then $b'$ is injective, and for $k<l$ we have
$(d-1)l-(d-1)k\ge d-1$. Hence $(b'(k)+t_1,b'(l)+t_d)\in \tilde A$ for all $k<l$, which
contradicts \cref{G:lem:Ztransfer}.
\end{proof}

To summarise: the configurations~\eqref{G:eq:corner} and~\eqref{G:eq:cornerd} cannot be
forced by the fibre-density hypotheses \eqref{G:eq:23}, \eqref{G:eq:I} or \eqref{G:eq:II}.
For $G=\Z$ (and for $G=\N$) this remains true after the lower corner points are dropped.
Whether a uniform hypothesis on the row fibres, for example syndeticity with a bound on the
gaps that does not depend on the fibre, suffices is open even for $G=\N$; see
\cref{G:rem:q510}(4). So is repair~(b) of \cref{G:rem:formulation}, which also requires
the columns to have positive density.
